A canonical concern about unilateral carbon pricing is \emph{leakage}: regulated firms raise prices, downstream buyers reallocate purchases toward less-regulated alternatives, and the net emissions reduction shrinks. The mechanism has two empirical preconditions. First, carbon costs must \emph{pass through} into the prices buyers pay. Second, conditional on pass-through, buyers must \emph{reallocate} their input mix in response to the relative-price change. Recent work has documented near-complete pass-through into the producer prices of regulated EU ETS sectors during Phase~III \citep{colmeretal2023,kanzig2023unequal}. The empirical question for the second precondition --- buyer reallocation --- is much less settled, particularly inside an integrated production network where reallocation can occur both across suppliers within a narrow input category and across categories within a buyer's input bill.

This paper provides a direct test of buyer-side reallocation using Belgian firm-to-firm transaction records. The Belgian setting is well-suited to this question for three reasons. First, the National Bank of Belgium maintains a near-universal panel of business-to-business invoices linked by VAT identifier; we observe the full domestic supplier composition of essentially every Belgian buyer. Second, ETS-regulated installations are a small but identifiable subset of suppliers ($\sim$280 firms, concentrated in a handful of NACE 2-digit sectors) for which we can measure pre-shock carbon-cost intensity using EUTL allowance and emissions records merged with Annual Accounts revenues. Third, the EU ETS supplies a pre-specified policy event with a well-understood discontinuity in policy stringency: the 2017 Market Stability Reserve reform triggered a roughly tenfold rise in EUA prices over 2017--2022, generating a Phase~IV cost shock with order-of-magnitude variation in pair-level pass-through exposure.

We examine reallocation at two margins. The first --- within a narrowly defined supplier sector (NACE 4-digit), \emph{across suppliers} --- corresponds to the substitution a buyer can make without changing the composition of its input basket. The second --- \emph{across} NACE 4-digit input categories within the same buyer-year --- corresponds to the substitution a buyer can make by reweighting categories of inputs (\eg substituting away from cement toward steel) when the average price in one category rises relative to others.

\paragraph{Findings.} Both margins yield null point estimates under the specifications closest to current best practice for difference-in-differences with continuous treatment and high-dimensional fixed effects:

\begin{itemize}
    \item \textbf{Within-NACE-4d, across-supplier substitution.} In the larger sample of 32,631 buyer $\times$ NACE-4d cells with at least one ETS supplier, the share of the most-carbon-exposed ETS supplier (the supplier with the largest pre-shock cost-share intensity in that 4-digit cell) does not change differentially with that supplier's exposure intensity after the 2015 cut-off. The pooled coefficient is $\beta = +1.34$ percentage points per unit of cost-share intensity (s.e.\ $0.96$, $p = 0.16$); after post-hoc detrending against a 2010--2014 linear pre-trend the coefficient falls to $+0.20$ (s.e.\ $1.14$, $p = 0.86$). In the most economically favorable subset --- cells where the supplier's NACE 4-digit accounts for the largest share of the buyer's total input bill, \ie where any per-unit price change would mechanically produce the largest dollar shock --- the coefficient is $\beta = +0.03$ (s.e.\ $1.39$, $p = 0.98$).

    \item \textbf{Across-NACE-4d, across-category substitution.} On a buyer $\times$ NACE 4-digit $\times$ year panel of 71.7 million cell-years, our preferred specification --- a binary regulated-category $\times$ \textit{Post} dummy with buyer $\times$ year and buyer $\times$ NACE 4-digit fixed effects, the most direct domestic analog of the regulated/unregulated event study in \citet{colmeretal2023} --- yields $\beta = -0.003$ (s.e.\ $0.009$, $p = 0.76$). The detrended event study under a continuous treatment shows uniformly \emph{positive} year coefficients in every clean post-2015 year (2015--2019 and 2021), pointing in the opposite direction from substitution.

    \item \textbf{Descriptive complement.} A standard Grossman--Krueger \citep{grossman1993} emissions decomposition on the same Belgian ETS firm panel attributes the 46\% aggregate emissions decline over 2005--2022 to scale, between-sector reallocation, within-sector reallocation, and within-firm abatement. Within-sector reallocation is essentially zero across all years and across both fixed-base and year-on-year specifications. This is consistent with --- though not equivalent to --- the buyer-side null at the within-NACE-4d margin.
\end{itemize}

\paragraph{What we do and do not claim.} The contribution of this paper is to document an observational null on a specific margin, not to estimate a structural parameter. We do not claim to identify the elasticity of substitution between regulated and unregulated suppliers; we do not claim that buyers \textit{would not} reallocate at higher carbon prices; and we do not claim a mechanism. Three caveats deserve emphasis up front:

\begin{enumerate}
    \item \emph{The within-NACE-4d test has a parallel-trends problem.} The event study around the 2015 cutoff shows significantly negative coefficients in 2010 ($\beta_{2010} = -5.70$, $p < 0.05$) and 2011 ($\beta_{2011} = -5.05$, $p < 0.05$): high-exposure suppliers were already losing share before the 2015 shock. The headline pooled coefficient on a smaller selected subsample (cells where the supplier was active throughout 2005--2022, 705 cells) is significant and negative ($\beta = -293$, $p = 0.03$) but is contaminated by this pre-trend. Our preferred specification is the larger 32{,}631-cell sample where the pooled coefficient is null and the post-hoc detrended coefficient is null, and we interpret the smaller-sample point estimate as not separately identified from the pre-trend.

    \item \emph{There is a small but statistically significant extensive-margin response.} An ancillary linear probability model on whether the most-exposed supplier remains active in the buyer-cell post-2015 yields $\beta = -3.61$ (s.e.\ $1.47$, $p = 0.014$). At typical exposure (90th percentile of cost-share intensity $\approx 0.0014$), the predicted probability of the supplier exiting the relationship rises by roughly half a percentage point. The point estimate is statistically detectable; the implied magnitude is small relative to baseline relationship-exit rates.

    \item \emph{Power is limited at the lower tail; the right tail is informative.} The distribution of pair-level cost shocks is heavy-tailed. The population p99 of the pair-shock-total measure (the percentage of buyer total-input-cost at risk under full pass-through) is $\approx 39\%$, comparable to or larger than the per-HS-product tariff cuts that \citet{peter2023} use to identify the elasticity of substitution. The tests reported here have economic power against \emph{large} substitution at the right tail and rule it out; they have less power against \emph{modest} substitution dispersed across the population.
\end{enumerate}

\paragraph{Related literature.} We connect to three strands. First, on EU ETS impacts on regulated firms: \citet{martin2014}, \citet{dechezlepretreetal2023}, \citet{colmeretal2023}, and \citet{kanzig2023unequal}. Our null at the buyer-side reallocation margin is consistent with the consistent finding in this literature that ETS firms reduced emissions without contracting output, which would be hard to reconcile with large reallocation toward unregulated competitors. Second, on input-output reallocation under cost shocks: \citet{huneeus2018} and \citet{peter2023}. Our null is, ex ante, surprising relative to the substitution magnitudes those papers identify; we discuss in Section~\ref{sec:discussion} the alignment of right-tail magnitudes between our setting and theirs. Third, on the firm-level dynamics of regulated supplier-customer relationships, we relate to the network-margin discussion in \citet{baqaee2019}.

The remainder of the paper proceeds as follows. Section~\ref{sec:data} describes the data and sample construction. Section~\ref{sec:decomposition} presents the descriptive supplier-side decomposition. Section~\ref{sec:test_h} reports the within-NACE-4d test. Section~\ref{sec:test_i} reports the across-NACE-4d test. Section~\ref{sec:discussion} bounds what the joint null can and cannot say. Section~\ref{sec:conclusion} concludes.
